% id: 13-01
% book: 베이직쎈 미적분1 · 01 함수의 극한
% section: 유형 001 함수의 수렴과 발산
% figure: none
% answer: 
% answer_source: 
% variant_level: 0 (원본 전사)
% 이 파일은 build.mjs 가 items.json 에서 만든다. 고칠 때는 items.json 을 고친다.
\smash{\raisebox{1.5mm}{\dmkichul{베이직쎈 미적분1 13쪽 01번}}}
$x\to 0$일 때, 함수 $f(x)$가 수렴하는 것만을 \textbf{보기}에서 있는 대로 고른 것은? \bogi{\begin{minipage}[t]{0.48\linewidth}\raggedright ㄱ.\hspace{1.5mm}% 13-01(13쪽) — 보기 ㄱ 의 그래프(원점을 지나 오른쪽 아래로 내려가는 직선). 스캔 화질이 낮아 TikZ 로 다시 그림.
% 원문에 있는 것만: 직선 하나 · 라벨 O, x, y, y=f(x). 눈금·좌표값은 원문에 없으므로 넣지 않는다.
\begin{tikzpicture}[baseline=(current bounding box.north), x=0.38cm, y=0.38cm, line width=0.55pt]
  \draw[->, >=stealth, line width=0.7pt] (-2.5,0) -- (2.6,0) node[below=1pt] {$x$};
  \draw[->, >=stealth, line width=0.7pt] (0,-2.5) -- (0,2.4) node[left=1pt] {$y$};
  \node[below left=-2pt and -3pt, font=\small] at (0,0) {$\pt{O}$};
  \draw (-1.95,1.95) -- (2.05,-1.85);
  \node[anchor=west, font=\small] at (0.2,-2.2) {$y=f(x)$};
\end{tikzpicture}
\end{minipage}\begin{minipage}[t]{0.48\linewidth}\raggedright ㄴ.\hspace{1.5mm}% 13-01(13쪽) — 보기 ㄴ 의 그래프(꼭짓점이 $(0,\,-2)$ 인 V 자 그래프). 스캔 화질이 낮아 TikZ 로 다시 그림.
% 원문에 있는 것만: y축 눈금 -2 · 두 반직선 · 라벨 O, x, y, y=f(x).
\begin{tikzpicture}[baseline=(current bounding box.north), x=0.38cm, y=0.38cm, line width=0.55pt]
  \draw[->, >=stealth, line width=0.7pt] (-2.55,0) -- (2.5,0) node[below=1pt] {$x$};
  \draw[->, >=stealth, line width=0.7pt] (0,-2.4) -- (0,2.6) node[left=1pt] {$y$};
  \draw (-1.5,1.75) -- (0,-2) -- (1.55,1.9);
  \node[above left=-3pt and -3pt, font=\small] at (0,0) {$\pt{O}$};
  \node[left=2pt, font=\small] at (0,-2) {$-2$};
  \node[anchor=west, font=\small] at (0.45,2.35) {$y=f(x)$};
\end{tikzpicture}
\end{minipage}\par\medskip\begin{minipage}[t]{0.48\linewidth}\raggedright ㄷ.\hspace{1.5mm}% 13-01(13쪽) — 보기 ㄷ 의 그래프(x축·y축을 점근선으로 하는 두 갈래 곡선 · 왼쪽 갈래는 x축 위, 오른쪽 갈래는 x축 아래).
% 스캔 화질이 낮아 TikZ 로 다시 그림. 원문에 있는 것만: 두 갈래 곡선 · 라벨 O, x, y, y=f(x)(눈금 없음).
\begin{tikzpicture}[baseline=(current bounding box.north), x=0.38cm, y=0.38cm, line width=0.55pt]
  \draw[->, >=stealth, line width=0.7pt] (-2.75,0) -- (2.65,0);
  \draw[->, >=stealth, line width=0.7pt] (0,-2.9) -- (0,2.65) node[right=1pt] {$y$};
  \draw[line width=0.6pt, domain=-2.7:-0.21, samples=80, smooth] plot (\x, {-0.5/\x});
  \draw[line width=0.6pt, domain=0.21:2.7, samples=80, smooth] plot (\x, {-0.5/\x});
  \node[above=1pt, font=\small] at (2.25,0) {$x$};
  \node[below left=-3pt and -3pt, font=\small] at (0,0) {$\pt{O}$};
  \node[anchor=west, font=\small] at (0.25,-2.3) {$y=f(x)$};
\end{tikzpicture}
\end{minipage}\begin{minipage}[t]{0.48\linewidth}\raggedright ㄹ.\hspace{1.5mm}% 13-01(13쪽) — 보기 ㄹ 의 그래프(원점에서 채운 점으로 끊기고 $(0,\,1)$ 의 빈 점에서 다시 시작하는 두 반직선).
% 스캔 화질이 낮아 TikZ 로 다시 그림. 원문의 빈 점(○)·채운 점(●) 과 y축 눈금 1 을 그대로 옮긴다.
\begin{tikzpicture}[baseline=(current bounding box.north), x=0.28cm, y=0.28cm, line width=0.55pt]
  \draw[->, >=stealth, line width=0.7pt] (-3.4,0) -- (3.3,0) node[below=1pt] {$x$};
  \draw[->, >=stealth, line width=0.7pt] (0,-2.75) -- (0,3.75) node[left=1pt] {$y$};
  \draw (-2.9,-2.4) -- (0,0);
  \draw (0,1) -- (2.4,3.0);
  \fill (0,0) circle (2.2pt);
  \draw[fill=white] (0,1) circle (2.2pt);
  \node[left=1pt, font=\small] at (0,1) {$1$};
  \node[below right=-2pt and -2pt, font=\small] at (0,0) {$\pt{O}$};
  \node[anchor=west, font=\small] at (0.5,3.6) {$y=f(x)$};
\end{tikzpicture}
\end{minipage}}
\dmchoicesauto{}{ㄱ, ㄴ}{ㄱ, ㄷ}{ㄴ, ㄹ}{ㄱ, ㄴ, ㄷ}{ㄱ, ㄴ, ㄹ}
